% GENERATED FILE. Edit canonical lessons, then run npm run generate:books.
% canonical-source-hash: fnv1a64:c5e3dd71a570a4af
% canonical-lessons: FR-C402-pare-brise, FR-C402-retroviseur, FR-C402-klaxon, FR-C402-klaxonner, FR-C402-embuer

\chapter{Bien voir la route}
\label{ch:fr-bien-voir-la-route}
\hlchaptermodality{402}

\begin{chapteropening}
\textbf{By the end of this chapter:} \emph{I can talk about seeing clearly from a car: the windscreen, the rear-view mirror, the horn, honking and windows that mist up.}

Complete the last lesson of chapter 402: I can talk about seeing clearly from a car: the windscreen, the rear-view mirror, the horn, honking and windows that mist up.
\end{chapteropening}

\section[le pare-brise]{le pare-brise — the windscreen}

\label{lesson:FR-C402-pare-brise}



\begin{quote}
Before the new one: say the French for to stall, then the French for the pedal.
\end{quote}

\subsection*{You'll want to know: le pare-brise}
\textbf{le pare-brise} — \textquotedblleft{}the windscreen\textquotedblright{}.

\begin{quote}
Un caillou a cassé le pare-brise.
\end{quote}

\begin{quote}
\emph{A stone broke the windscreen.}
\end{quote}

The plural does not change: \emph{des pare-brise}.

\begin{cousinweb}
\emph{pare-brise} is built from \emph{parer} 'to ward off' and \emph{brise} 'breeze' (EtymDB): the glass that keeps the wind out. \emph{Parer} goes back to Latin \emph{parāre} 'to prepare' (Lewis \& Short).
\end{cousinweb}

\begin{grammarlens}[title={What you've built}]
One more word for telling, explaining and describing at B1.
\end{grammarlens}

\subsection*{Your turn}
\begin{itemize}
\raggedright
  \item \emph{Say it:} \emph{le pare-brise}
  \item \emph{Say it:} \emph{le pare-brise}, once more
  \item \emph{Say it:} a sentence of your own with \emph{le pare-brise}
  \item \emph{Recall:} say the French for to stall, then the French for the pedal, then say \emph{le pare-brise} again
\end{itemize}

\begin{tcolorbox}[breakable,colback=teal!4,colframe=teal!35!black,title={Before you move on}]
What does le pare-brise mean? (\textquotedblleft{}The windscreen\textquotedblright{}.) Say it once more.
\end{tcolorbox}
\section[le rétroviseur]{le rétroviseur — the rear-view mirror; the wing mirror}

\label{lesson:FR-C402-retroviseur}



\begin{quote}
Before the new one: say the French for the pedal, then the French for the windscreen.
\end{quote}

\subsection*{You'll want to know: le rétroviseur}
\textbf{le rétroviseur} — \textquotedblleft{}the rear-view mirror; the wing mirror\textquotedblright{}.

\begin{quote}
Avant de tourner, il regarde dans le rétroviseur.
\end{quote}

\begin{quote}
\emph{Before turning, he checks the mirror.}
\end{quote}

In speech, often \emph{le rétro}.

\begin{cousinweb}
\emph{rétroviseur} is built from \emph{rétro-}, from Latin \emph{retrō} 'backwards', and \emph{viseur} 'sight', from \emph{viser} 'to aim', which goes back to Latin \emph{vidēre} 'to see' (EtymDB).
\end{cousinweb}

\begin{grammarlens}[title={What you've built}]
One more word for telling, explaining and describing at B1.
\end{grammarlens}

\subsection*{Your turn}
\begin{itemize}
\raggedright
  \item \emph{Say it:} \emph{le rétroviseur}
  \item \emph{Say it:} \emph{le rétroviseur}, once more
  \item \emph{Say it:} a sentence of your own with \emph{le rétroviseur}
  \item \emph{Recall:} say the French for the pedal, then the French for the windscreen, then say \emph{le rétroviseur} again
\end{itemize}

\begin{tcolorbox}[breakable,colback=teal!4,colframe=teal!35!black,title={Before you move on}]
What does le rétroviseur mean? (\textquotedblleft{}The rear-view mirror; the wing mirror\textquotedblright{}.) Say it once more.
\end{tcolorbox}
\section[le klaxon]{le klaxon — the horn (of a car)}

\label{lesson:FR-C402-klaxon}



\begin{quote}
Before the new one: say the French for the windscreen, then the French for the rear-view mirror.
\end{quote}

\subsection*{You'll want to know: le klaxon}
\textbf{le klaxon} — \textquotedblleft{}the horn (of a car)\textquotedblright{}.

\begin{quote}
Dans les rues étroites, on entend souvent le klaxon des taxis.
\end{quote}

\begin{quote}
\emph{In the narrow streets you often hear the taxis' horns.}
\end{quote}

First a trade name; now an ordinary noun.

\begin{cousinweb}
\emph{klaxon} was borrowed from English \emph{klaxon} (EtymDB), first the trade name of an electric car horn.
\end{cousinweb}

\begin{grammarlens}[title={What you've built}]
One more word for telling, explaining and describing at B1.
\end{grammarlens}

\subsection*{Your turn}
\begin{itemize}
\raggedright
  \item \emph{Say it:} \emph{le klaxon}
  \item \emph{Say it:} \emph{le klaxon}, once more
  \item \emph{Say it:} a sentence of your own with \emph{le klaxon}
  \item \emph{Recall:} say the French for the windscreen, then the French for the rear-view mirror, then say \emph{le klaxon} again
\end{itemize}

\begin{tcolorbox}[breakable,colback=teal!4,colframe=teal!35!black,title={Before you move on}]
What does le klaxon mean? (\textquotedblleft{}The horn (of a car)\textquotedblright{}.) Say it once more.
\end{tcolorbox}
\section[klaxonner]{klaxonner — to honk, to sound the horn}

\label{lesson:FR-C402-klaxonner}



\begin{quote}
Before the new one: say the French for the rear-view mirror, then the French for the horn.
\end{quote}

\subsection*{You'll want to know: klaxonner}
\textbf{klaxonner} — \textquotedblleft{}to honk, to sound the horn\textquotedblright{}.

\begin{quote}
Le chauffeur a klaxonné devant la maison.
\end{quote}

\begin{quote}
\emph{The driver honked outside the house.}
\end{quote}

Regular \emph{-er} verb: \emph{je klaxonne}.

\begin{cousinweb}
\emph{klaxonner} is built on \emph{klaxon} 'car horn' (EtymDB), which was borrowed from English \emph{klaxon}.
\end{cousinweb}

\begin{grammarlens}[title={What you've built}]
One more word for telling, explaining and describing at B1.
\end{grammarlens}

\subsection*{Your turn}
\begin{itemize}
\raggedright
  \item \emph{Say it:} \emph{klaxonner}
  \item \emph{Say it:} \emph{klaxonner}, once more
  \item \emph{Say it:} a sentence of your own with \emph{klaxonner}
  \item \emph{Recall:} say the French for the rear-view mirror, then the French for the horn, then say \emph{klaxonner} again
\end{itemize}

\begin{tcolorbox}[breakable,colback=teal!4,colframe=teal!35!black,title={Before you move on}]
What does klaxonner mean? (\textquotedblleft{}To honk, to sound the horn\textquotedblright{}.) Say it once more.
\end{tcolorbox}
\section[embuer]{embuer — to mist up (mostly reflexive: s’embuer)}

\label{lesson:FR-C402-embuer}



\begin{quote}
Before the new one: say the French for the horn, then the French for to honk.
\end{quote}

\subsection*{You'll want to know: embuer}
\textbf{embuer} — \textquotedblleft{}to mist up (mostly reflexive: s’embuer)\textquotedblright{}.

\begin{quote}
Quand il fait froid, les vitres s'embuent.
\end{quote}

\begin{quote}
\emph{When it is cold, the windows mist up.}
\end{quote}

Mostly reflexive: \emph{les vitres s'embuent}. \emph{Embué} — misty.

\begin{cousinweb}
\emph{embuer} is built on \emph{buée} 'mist, condensation'. \emph{Buée} is usually traced to an old verb \emph{buer} 'to do the washing', of Germanic origin: the steam that rose from the washing tubs. EtymDB does not list the word.
\end{cousinweb}

\begin{grammarlens}[title={What you've built}]
One more word for telling, explaining and describing at B1.
\end{grammarlens}

\subsection*{Your turn}
\begin{itemize}
\raggedright
  \item \emph{Say it:} \emph{embuer}
  \item \emph{Say it:} \emph{embuer}, once more
  \item \emph{Say it:} a sentence of your own with \emph{embuer}
  \item \emph{Recall:} say the French for the horn, then the French for to honk, then say \emph{embuer} again
\end{itemize}

\begin{tcolorbox}[breakable,colback=teal!4,colframe=teal!35!black,title={Before you move on}]
What does embuer mean? (\textquotedblleft{}To mist up (mostly reflexive: s’embuer)\textquotedblright{}.) Say it once more.
\end{tcolorbox}
